\section{Introduction}
\label{sec:introduction}

Temporal computation requires selective memory: recent inputs must remain
available long enough to be useful, while distant inputs and the processor's
starting state must fade. Reservoir computing achieves this balance with fixed
dynamics and trains only a linear map from measured signals to the desired
output~\cite{Maass2002,JaegerESP2001}. Quantum reservoir computing (QRC) uses a
quantum system to transform each input into measured features~\cite{Fujii2017}.
Because the internal dynamics are not trained, how they retain input history is
central to the computation. Purely unitary evolution does not by itself provide
the required forgetting. Coupling to an environment can therefore be functional
rather than merely detrimental: dissipative dynamics can implement nonlinear
transformations whose dependence on older inputs gradually
vanishes~\cite{ChenNurdin2019}.

Engineered dissipation already prepares and stabilizes states, protects quantum
memories, and drives computation
~\cite{Poyatos1996,Verstraete2009,Pastawski2011,Harrington2022}. In QRC, noise,
damping, optical absorption, measurement backaction, feedback, and
reinitialization can create or regulate temporal memory
~\cite{Chen2020,Sannia2024,Ricci2026,Gotting2025,Mujal2023,Franceschetto2026,
KobayashiFeedback2024,Cindrak2024,SanniaNonMarkov2026,Paparelle2026}. Studies of
coherence and measurements made during a circuit further clarify practically
accessible regimes~\cite{Palacios2024,Hu2024}. Sannia \emph{et al.}, for example, showed
that tuning uniform, independent qubit relaxation can improve a continuously
driven spin reservoir~\cite{Sannia2024}. Related work likewise tunes
scalar controls such as drive strength, interaction strength, input
duration, and damping rate~\cite{MartinezPenaPRL2021,Baumann2026}. These studies
establish that dissipation can help, but they largely choose the environmental
process in advance and ask how strongly it should act.

An environment specifies more than a rate. It also determines which system
transitions lose information separately and which couple through the same
channel. With uniform local relaxation, each qubit has its own route into the
environment. With collective relaxation, one shared channel acts directly on a collective
combination of qubit transitions. Other combinations reach it only after the
Hamiltonian mixes them together. This distinction is familiar from collective
emission and structures protected from noise in quantum
networks~\cite{Dicke1954,Cabot2018}. Its computational consequence for a
quantum system driven by inputs is less clear. Does the environment merely set how
quickly a processor forgets, or can its pattern of coupling select which parts
of an input history remain available?

Here we answer this question numerically for finite reservoirs of driven spins. We
show that changing the organization of environmental coupling changes what
earlier inputs a fixed readout can recover. We vary which transitions share a
channel and how a common, explicitly defined coupling budget is distributed
among them. Within each principal pair, the Hamiltonian, input sequence,
measured Pauli observables, training procedure, and data split are identical.
Across six types of environmental process, no design dominates every benchmark.
Instead, different couplings favor memory, nonlinear processing, parity, and
chaotic forecasting in different ways. The environment therefore changes the
kind of temporal information available to the readout, not simply whether one
score improves.

The clearest recurring contrast is between collective and uniform local relaxation.
Collective relaxation leaves more recent input history accessible to the linear
readout and improves recovery of a nonlinear function of that history. Both
effects appear in every principal paired instance with five qubits. The ordering
remains after the longer washout, when tested starting-state effects are
negligible relative to the task contrast.
It recurs from four to eight qubits, across four spin Hamiltonian ensembles, and
after replacing continuous input driving with encoding through reset operations. It also
survives separate comparisons that match average dissipative activity or a
representative relaxation timescale, and when each design independently selects
its operating point from a bounded range. No obvious scalar measure of
dissipation therefore explains the result.

The most direct test changes only the collective transition addressed by a
shared channel. Rotating this direction changes accessible memory even though
the number and nonzero strengths of the coupled directions, the total weight,
and the weight on every qubit all remain fixed. The dependence on coupling
direction recurs in a separate study with six qubits. Redistributing coupling weight
within one type of process likewise changes memory. Thus even detailed coarse
descriptions of dissipation do not fix the computation: which collective
transition the environment directly addresses is itself a design variable. The
conclusions concern the tested finite reservoirs. They do not establish a
universal optimum, an asymptotic scaling law, or a unique microscopic cause. A
common coupling budget is a controlled comparison, not complete physical
equivalence.

These findings connect QRC with the broader practice of engineering
couplings between systems and environments for state preparation, information processing,
and memory protection
~\cite{Poyatos1996,Verstraete2009,Pastawski2011,Harrington2022,Cattaneo2019,
Brehm2021,Sheremet2023,CattaneoPRX2023}. Here, memory means information about
earlier inputs that remains linearly recoverable from the specified Pauli
observables after the settling period. It is not a claim about quantum memory
lifetime. Within this scope, the architectural implication is clear. The
environment is not merely an error source or a clock that sets how quickly
information disappears. The transitions it can access shape which parts of the
past remain usable. Environmental coupling should therefore be designed
alongside the Hamiltonian, input encoding, and measurement. For temporal quantum
processing, the environment is not merely the boundary of the processor. It is
part of the processor.
